\documentclass[11pt]{amsart}
\usepackage{amsmath,amssymb,amsthm}
\usepackage{mathrsfs}
\usepackage{hyperref}
\usepackage{geometry}
\usepackage{booktabs}
\geometry{margin=1in}

%--- Theorem environments
\newtheorem{theorem}{Theorem}[section]
\newtheorem{proposition}[theorem]{Proposition}
\newtheorem{lemma}[theorem]{Lemma}
\newtheorem{corollary}[theorem]{Corollary}
\newtheorem{definition}[theorem]{Definition}
\newtheorem{remark}[theorem]{Remark}
\newtheorem{conjecture}[theorem]{Conjecture}

%--- Operators
\renewcommand{\Re}{\operatorname{Re}}
\renewcommand{\Im}{\operatorname{Im}}
\newcommand{\LP}{\mathrm{LP}}
\newcommand{\diag}{\operatorname{diag}}

\title[A Curvature Criterion for RH]{A curvature criterion for the Riemann Hypothesis\\
  and the unification of four analytic frameworks}
\author{David Alejandro Trejo Pizzo}
\thanks{Independent Researcher. \texttt{dtrejopizzo@gmail.com}}


\begin{document}

\begin{abstract}
We derive an explicit curvature criterion equivalent to the Riemann Hypothesis (RH) and
use it to unify four analytic frameworks that have separately been proposed as proof
strategies: the de~Bruijn--Newman heat flow and Kre\u{\i}n monotonicity~(A), the
APS spectral-flow approach~(B), the $K$-theoretic obstruction class~(C), and a
transversal energy functional~(D).

The central new result is the \emph{curvature formula}: for $s = \sigma + i\gamma$
with $\sigma = 1/2$ and $\gamma \notin \{\gamma_\rho\}$,
\[
  \mathcal{C}(\gamma)
  \;:=\;
  \frac{\partial^2}{\partial\sigma^2}\log|\xi(\sigma+i\gamma)|
  \bigg|_{\sigma=1/2}
  \;=\;
  \sum_\rho \frac{(\gamma-\gamma_\rho)^2 - (\sigma - \beta_\rho)^2}{|s-\rho|^4},
\]
where the sum runs over all non-trivial zeros $\rho = \beta_\rho + i\gamma_\rho$ of~$\zeta$.
We prove:
\[
  \mathrm{RH}
  \;\iff\;
  \mathcal{C}(\gamma) > 0 \quad \forall\,\gamma \notin \{\gamma_\rho\},
\]
or equivalently, $f(t) := \xi(1/2+it)$ belongs to the Laguerre--P\'olya class~$\LP$.

We further prove the explicit sector decomposition
$G(t) = G_{\mathrm{poly}}(t) + G_\Gamma(t) + G_\zeta(t)$ of the function
$G(t) = (\xi'/\xi)'(1/2+it)$, with exact formulas for the analytic sector
($G_{\mathrm{poly}}$ and $G_\Gamma$ computable from the functional equation alone)
and identify the arithmetic sector $G_\zeta$ as the irreducible obstruction.
The paper concludes with a precise statement of the remaining wall: proving
$G_\zeta(t) > -G_{\mathrm{ana}}(t)$ for all $t > t_*$ without assuming RH
is equivalent to RH.
\end{abstract}

\maketitle

\tableofcontents
\newpage

%------------------------------------------------------------------
\section{Introduction}
\label{sec:intro}
%------------------------------------------------------------------

The Riemann Hypothesis asserts that every non-trivial zero $\rho$ of the Riemann
$\zeta$-function satisfies $\Re(\rho) = 1/2$.
Since Riemann's memoir of 1859, this statement has accumulated more than a dozen
equivalent reformulations, each exposing a different face of the same problem.
The present paper contributes a new face---the \emph{curvature diagnostic}---and
shows that four independent frameworks proposed in the literature
as potential proof strategies are all equivalent to this single analytic criterion.

\subsection*{Main contributions}

\begin{enumerate}
  \item \textbf{Curvature formula (Theorem~\ref{thm:curvature}).} An explicit formula
    for $\partial^2_{\sigma\sigma}\log|\xi(\sigma+i\gamma)|$ in terms of the non-trivial zeros.
  \item \textbf{Exact equivalence (Theorem~\ref{thm:rh-curvature}).} $\mathrm{RH}
    \iff \mathcal{C}(\gamma) > 0$ for all $\gamma \notin \{\gamma_\rho\}$.
  \item \textbf{LP-class unification (Theorem~\ref{thm:lp}).} $f(t) = \xi(1/2+it)$
    belongs to the Laguerre--P\'olya class $\iff$ $\mathcal{C}(t) > 0$ $\iff$ all
    four analytic frameworks (A)--(D) hold simultaneously.
  \item \textbf{Sector decomposition (Theorem~\ref{thm:sectors}).} Exact closed-form
    expressions for $G_{\mathrm{poly}}(t)$ and $G_\Gamma(t)$; precise identification
    of $G_\zeta(t)$ as the irreducible arithmetic obstruction.
  \item \textbf{Negative result on $\sigma$-concavity (Proposition~\ref{prop:sigma-convex}).}
    Under RH, $\log|\xi(\sigma+i\gamma)|$ is $\sigma$-\emph{convex} (not concave) on the
    critical line; the correct direction is $t$-concavity.
\end{enumerate}

\subsection*{Relation to known results}

The equivalence $\mathrm{RH} \iff \xi(1/2+it) \in \LP$ is due to de~Bruijn~\cite{deBruijn1950}
(see also Newman~\cite{Newman1976}). The inequality $\Lambda \geq 0$ was proved by
Rodgers--Tao~\cite{RodgersTao2020}. The connection between spectral flow and $\kappa(Q)$
uses the Atiyah--Patodi--Singer theorem~\cite{APS1975}.
The new content of the present paper is:
(i) the explicit Hadamard-product formula for $\mathcal{C}(\gamma)$ and its sign analysis;
(ii) the resolution of the direction of concavity (Proposition~\ref{prop:sigma-convex});
(iii) the exact sector decomposition with polygamma formulas;
(iv) the synthesis showing A--D reduce to the same LP-class problem.

\subsection*{Notation}
Throughout, $\xi(s) = \frac{1}{2}s(s-1)\pi^{-s/2}\Gamma(s/2)\zeta(s)$ is the completed
zeta function; non-trivial zeros of $\zeta$ are written $\rho = \beta_\rho + i\gamma_\rho$
with $0 < \beta_\rho < 1$; $\psi^{(1)}$ denotes the trigamma function.
We write $f(t) := \xi(1/2+it)$ for $t \in \mathbb{R}$.

%------------------------------------------------------------------
\section{The Hadamard curvature formula}
\label{sec:curvature}
%------------------------------------------------------------------

\subsection{Setup}

The function $\xi$ is entire of order one, with Hadamard product
\begin{equation}\label{eq:hadamard}
  \xi(s) = \xi(0)\prod_\rho\!\left(1-\frac{s}{\rho}\right)
  \quad\text{(convergent in pairs $\rho,1-\rho$)}.
\end{equation}
Write $\mathcal{L}(s) = \log|\xi(s)|$ (harmonic away from zeros).
By the Hadamard product,
\begin{equation}\label{eq:log-xi}
  \mathcal{L}(s) = \log|\xi(0)| + \sum_\rho \Re\log\!\left(1-\frac{s}{\rho}\right)+\frac{s}{\rho}.
\end{equation}

\begin{theorem}[Curvature formula]\label{thm:curvature}
For $s = \sigma + i\gamma$ with $\sigma \neq \beta_\rho$ for all $\rho$,
\begin{equation}\label{eq:curvature-cartesian}
  \frac{\partial^2}{\partial\sigma^2}\,\mathcal{L}(\sigma+i\gamma)
  \;=\;
  \sum_\rho \frac{(\gamma-\gamma_\rho)^2 - (\sigma-\beta_\rho)^2}{|s-\rho|^4}.
\end{equation}
\end{theorem}

\begin{proof}
Differentiate~\eqref{eq:log-xi} twice in $\sigma$.  From a single factor:
\[
  \frac{\partial^2}{\partial\sigma^2}\Re\log\!\left(1-\frac{s}{\rho}\right)
  = \Re\frac{-1}{(s-\rho)^2}
  = -\frac{(\sigma-\beta_\rho)^2-(\gamma-\gamma_\rho)^2}{|s-\rho|^4}
  = \frac{(\gamma-\gamma_\rho)^2-(\sigma-\beta_\rho)^2}{|s-\rho|^4}.
\]
Summing over $\rho$ (the series converges absolutely for $s$ away from zeros
by the Riemann--von Mangoldt density $N(T) = O(T\log T)$) gives~\eqref{eq:curvature-cartesian}.
\end{proof}

\begin{remark}
The angular form is equivalent: writing $s - \rho = |s-\rho|e^{i\theta_\rho}$,
the numerator equals $|s-\rho|^2 \cos(2\theta_\rho)$, so
$\partial^2_{\sigma\sigma}\mathcal{L} = -\sum_\rho \cos(2\theta_\rho)/|s-\rho|^2$.
The sign of each term is determined by the angle $\theta_\rho$: positive in the
``vertical cone'' $|\theta_\rho| > \pi/4$, negative in the ``horizontal cone''
$|\theta_\rho| < \pi/4$.
\end{remark}

\subsection{Sign on the critical line under RH}

\begin{proposition}\label{prop:sigma-convex}
Assume RH. Then for all $\gamma \notin \{\gamma_\rho\}$,
\[
  \frac{\partial^2}{\partial\sigma^2}\,\mathcal{L}\!\left(\tfrac{1}{2}+i\gamma\right) > 0.
\]
That is, $\sigma \mapsto \log|\xi(\sigma+i\gamma)|$ is \emph{strictly convex} in $\sigma$
at $\sigma = 1/2$.
\end{proposition}

\begin{proof}
Under RH, $\beta_\rho = 1/2$ for all $\rho$, so $\sigma - \beta_\rho = 0$
at $\sigma = 1/2$.  Each term in~\eqref{eq:curvature-cartesian} becomes
$(\gamma-\gamma_\rho)^2/|s-\rho|^4 = 1/(\gamma-\gamma_\rho)^2 > 0$.
\end{proof}

\begin{remark}[Resolution of the ``wrong direction'' error]
A natural initial intuition is that RH should imply \emph{concavity} of
$\log|\xi|$ in $\sigma$ (since the critical line is a ``preferred'' locus).
Proposition~\ref{prop:sigma-convex} shows this is \emph{false}: RH implies
\emph{convexity}. The correct direction for LP-class characterization is
$t$-concavity, not $\sigma$-concavity (see Section~\ref{sec:lp}).
\end{remark}

%------------------------------------------------------------------
\section{The curvature criterion and the Laguerre--P\'olya class}
\label{sec:lp}
%------------------------------------------------------------------

\subsection{The function $f(t) = \xi(1/2+it)$ is real-valued}

\begin{proposition}\label{prop:real}
For all $t \in \mathbb{R}$, $\xi(1/2+it) \in \mathbb{R}$.
\end{proposition}

\begin{proof}
Combining the functional equation $\xi(s) = \xi(1-s)$ with the reality
$\overline{\xi(s)} = \xi(\bar{s})$ (entire function with real Taylor coefficients at
$s = 1/2$):
$\overline{\xi(1/2+it)} = \xi(1/2-it) = \xi(1-(1/2+it)) = \xi(1/2+it)$.
\end{proof}

\subsection{The Laguerre--P\'olya class}

\begin{definition}
An entire function $f: \mathbb{C} \to \mathbb{C}$ belongs to the
\emph{Laguerre--P\'olya class}~$\LP$ if it is a uniform limit on compact sets of
polynomials with only real zeros, or equivalently if it has the form
\[
  f(t) = c\,t^m e^{-\alpha t^2+\beta t}\prod_k\!\left(1-\frac{t}{t_k}\right)e^{t/t_k},
\]
with $c,\beta,t_k \in \mathbb{R}$, $\alpha \geq 0$, $\sum_k t_k^{-2} < \infty$.
\end{definition}

The class $\LP$ consists exactly of the entire real functions with all real zeros.
The following is the fundamental equivalence (de~Bruijn~\cite{deBruijn1950},
Newman~\cite{Newman1976}):

\begin{proposition}\label{prop:rh-lp}
The following are equivalent:
\begin{enumerate}
  \item[\rm(1)] RH: every non-trivial zero of $\zeta$ satisfies $\Re(\rho) = 1/2$.
  \item[\rm(2)] $f(t) = \xi(1/2+it)$ belongs to~$\LP$.
\end{enumerate}
\end{proposition}

\subsection{The curvature criterion}

\begin{theorem}[Curvature criterion]\label{thm:rh-curvature}
Define $\mathcal{C}(\gamma) := -\partial^2_{tt}\log|f(t)||_{t=\gamma}
= \partial^2_{\sigma\sigma}\mathcal{L}(1/2+i\gamma)$ (using harmonicity of $\mathcal{L}$).
Then:
\[
  f \in \LP
  \;\iff\;
  \mathcal{C}(\gamma) > 0
  \quad\forall\,\gamma \notin \{\gamma_\rho\}.
\]
In particular, $\mathrm{RH} \iff \mathcal{C}(\gamma) > 0$ for all such $\gamma$.
\end{theorem}

\begin{proof}
For $f \in \LP$, the Hadamard product gives all $t_k$ real, hence
$\partial^2_{tt}\log|f(t)| = -2\alpha - \sum_k (t-t_k)^{-2} < 0$,
i.e., $\mathcal{C}(t) > 0$.

Conversely, suppose $f$ has a non-real zero $t_0 = \gamma_0 + ib_0$ ($b_0 > 0$).
The conjugate pair $\{t_0, \bar t_0\}$ contributes
\[
  \Re\!\left(\frac{-1}{(t-t_0)^2} - \frac{1}{(t-\bar t_0)^2}\right)
  = \frac{-2(t-\gamma_0)^2 + 2b_0^2}{[(t-\gamma_0)^2+b_0^2]^2},
\]
which is \emph{positive} for $|t-\gamma_0| < b_0$.  Thus $\partial^2_{tt}\log|f| > 0$
(i.e., $\mathcal{C}(t) < 0$) in the strip $(\gamma_0-b_0, \gamma_0+b_0)$.
Hence $f \notin \LP$ implies $\mathcal{C} < 0$ somewhere.
\end{proof}

\begin{corollary}
$\mathrm{RH} \iff \mathcal{C}(t) = \displaystyle\sum_\rho \frac{1}{(t-\gamma_\rho)^2}$
for all $t \notin \{\gamma_\rho\}$.
\end{corollary}

\begin{proof}
Under RH, $\beta_\rho = 1/2$, so Theorem~\ref{thm:curvature} gives
$\mathcal{C}(t) = \sum_\rho 1/(t-\gamma_\rho)^2 > 0$. Conversely, if
$\mathcal{C}(t) > 0$ for all $t$, Theorem~\ref{thm:rh-curvature} gives $f \in \LP$,
hence RH by Proposition~\ref{prop:rh-lp}.
\end{proof}

%------------------------------------------------------------------
\section{The four analytic frameworks and their unification}
\label{sec:fronts}
%------------------------------------------------------------------

We now show that four frameworks separately proposed in the literature as potential
RH proof strategies reduce to a single condition: $f(t) \in \LP$.

\subsection{Framework A: de Bruijn--Newman flow and Kre\u{\i}n monotonicity}

The de~Bruijn--Newman constant $\Lambda \in \mathbb{R}$ is defined as the infimum of
$t \geq 0$ for which $H_t(z) = \int_0^\infty \Phi(u)e^{tu^2}\cos(uz)\,du$ has only
real zeros, where $\Phi$ is the Fourier transform weight of $\xi$.
The key results are:
\begin{itemize}
  \item $\mathrm{RH} \iff \Lambda = 0$ (de~Bruijn~\cite{deBruijn1950}, Newman~\cite{Newman1976}).
  \item $\Lambda \geq 0$ (Rodgers--Tao~\cite{RodgersTao2020}).
  \item The flow $t \mapsto H_t$ is the gradient flow of the log-repulsion energy
    $E_{\log} = -\sum_{j\neq k}\log|z_j-z_k|$ (Csordas--Norfolk--Varga~\cite{CNV1988}).
\end{itemize}
The critical line is the \emph{attractor} of the de~Bruijn--Newman flow.
Framework A reduces to: is $\Lambda = 0$?  This is equivalent to $H_0 = f \in \LP$.

\subsection{Framework B: APS spectral flow}

Let $\{T_\lambda\}_{\lambda \in [0,1]}$ be a family of self-adjoint operators on $L^2(\mathbb{R})$
associated to the Weil explicit formula, with $T_\lambda$ passing through zero eigenvalues.
The Atiyah--Patodi--Singer spectral flow $\mathrm{sf}(\{T_\lambda\})$ counts signed eigenvalue
crossings.

\begin{proposition}
$\mathrm{sf}(\{T_\lambda\}) = \kappa(Q)$, where $\kappa(Q)$ is the Kre\u{\i}n negative
index of the Weil quadratic form $Q$.
\end{proposition}

By the APS formula, $\kappa = \mathrm{Ind}(D_{\mathrm{APS}}) + \eta(T_0)/2$.
Framework B reduces to: is $\kappa = 0$?  Since $\kappa = 0 \iff$ no off-line zeros
$\iff$ RH $\iff f \in \LP$, Framework B reduces to the same LP condition.

\subsection{Framework C: $K$-theoretic obstruction}

The Connes scaling operator $H_C$ acts on $L^2(C_\mathbb{Q})$ where
$C_\mathbb{Q} = \mathbb{A}_\mathbb{Q}^*/\mathbb{Q}^*$ is the id\`ele class group.
The Baum--Connes isomorphism (which holds for $C_\mathbb{Q}$ since it is amenable) gives
$K_0(C^*(C_\mathbb{Q})) \cong \mathbb{Z}$.
The projection $P_- = \mathbf{1}_{(-\infty,0)}(Q)$ onto the negative part of $Q$ defines
an element $[P_-] \in K_0(C^*(C_\mathbb{Q}))$.

\begin{proposition}
$[P_-] = \kappa \in K_0 \cong \mathbb{Z}$.
Hence $[P_-] = 0 \iff \kappa = 0 \iff f \in \LP \iff \mathrm{RH}$.
\end{proposition}

Framework C reduces to: is $[P_-] = 0$ in $K_0$?

\subsection{Framework D: transversal energy functional}

Define the transversal energy $\Phi = \sum_j b_j^2$ (where $b_j = \Re(\rho_j) - 1/2$
for off-line zeros) and the arithmetic energy $E_\mathrm{arith} = E_\mathrm{log} + E_\mathrm{conf} + \Phi$.
By symmetry of the Weil form, $b_j = 0$ is always a \emph{stationary point} of
$E_\mathrm{arith}$.  The question is whether it is a \emph{minimum}.

The energy diagnostic is:
\[
  E_{\mathrm{curv}}(f) := \int_\mathbb{R} \mathcal{C}(t)\,dt.
\]

\begin{proposition}
$E_{\mathrm{curv}}(f) > 0$ for all $f \in \LP$.
Moreover, $\mathcal{C}(t) > 0$ for all $t$ $\iff$ $f \in \LP$ $\iff$ $\mathrm{RH}$.
\end{proposition}

Framework D reduces to: is the Hessian of $E_\mathrm{arith}$ in the transversal $b_j$-direction
positive definite?  The answer is yes $\iff f \in \LP \iff$ RH.

\subsection{The unification theorem}

\begin{theorem}[Unification]\label{thm:lp}
The following conditions are equivalent:
\begin{enumerate}
  \item[\rm(A)] $\Lambda = 0$ (de Bruijn--Newman).
  \item[\rm(B)] $\kappa(Q) = 0$ (spectral flow / Kre\u{\i}n index).
  \item[\rm(C)] $[P_-] = 0$ in $K_0(C^*(C_\mathbb{Q}))$.
  \item[\rm(D)] $\mathcal{C}(t) > 0$ for all $t \notin \{\gamma_\rho\}$ (curvature criterion).
  \item[\rm(E)] $f(t) = \xi(1/2+it) \in \LP$.
  \item[\rm(F)] RH.
\end{enumerate}
\end{theorem}

\begin{proof}
(F)$\iff$(E): Proposition~\ref{prop:rh-lp} (de~Bruijn--Newman).
(E)$\iff$(D): Theorem~\ref{thm:rh-curvature}.
(D)$\iff$(F): by the equivalences above.
(A)$\iff$(F): $\Lambda = 0 \iff H_0 \in \LP \iff$ (E).
(B)$\iff$(F): $\kappa = \#\{\text{off-line zero orbits}\}\cdot 2 = 0 \iff$ RH.
(C)$\iff$(F): Proposition above.
\end{proof}

%------------------------------------------------------------------
\section{The sector decomposition of $G(t)$}
\label{sec:sectors}
%------------------------------------------------------------------

Define $G(t) := (\xi'/\xi)'(1/2+it)$.  By the Hadamard product,
$G(t) = \sum_\rho d^2/dt^2 \log|1/2+it-\rho|$, and by harmonicity,
$G(t) = \mathcal{C}(t)$.  We decompose $G$ according to the three factors of $\xi$.

\begin{theorem}[Sector decomposition]\label{thm:sectors}
\[
  G(t) = G_{\mathrm{poly}}(t) + G_\Gamma(t) + G_\zeta(t),
\]
where:
\begin{align}
  G_{\mathrm{poly}}(t) &= \frac{2(1/4 - t^2)}{(t^2+1/4)^2}, \label{eq:gpoly}\\[4pt]
  G_\Gamma(t) &= -\tfrac{1}{4}\,\Re\!\left[\psi^{(1)}\!\left(\tfrac{1}{4}+\tfrac{it}{2}\right)\right],
  \label{eq:ggamma}\\[4pt]
  G_\zeta(t) &= \frac{d^2}{dt^2}\log|\zeta\!\left(\tfrac{1}{2}+it\right)|.
  \label{eq:gzeta}
\end{align}
\end{theorem}

\begin{proof}
From $\log|\xi(s)| = \log|\tfrac{1}{2}s(s-1)| - \tfrac{\log\pi}{2}\sigma + \log|\Gamma(s/2)| + \log|\zeta(s)|$,
differentiate twice in $t$ and use harmonicity ($\partial_{tt} = -\partial_{\sigma\sigma}$).

\medskip\noindent\textit{Polynomial sector.} $|\tfrac{1}{2}s(s-1)|_{s=1/2+it} = \tfrac{1}{2}(t^2+1/4)$, so
$\log|\tfrac{1}{2}s(s-1)| = \log\tfrac{1}{2}+\log(t^2+1/4)$, and~\eqref{eq:gpoly} follows.

\medskip\noindent\textit{Gamma sector.} $d^2/dt^2\log|\Gamma(1/4+it/2)| = \Re[(i/2)^2\psi^{(1)}(1/4+it/2)]
= -\tfrac{1}{4}\Re\psi^{(1)}(1/4+it/2)$, giving~\eqref{eq:ggamma}.
\end{proof}

\subsection{Properties of the analytic sector $G_{\mathrm{ana}} = G_{\mathrm{poly}} + G_\Gamma$}

\begin{proposition}\label{prop:gana}
The analytic sector satisfies:
\begin{enumerate}
  \item[\rm(i)] $G_{\mathrm{poly}}(t) > 0$ for $|t| < 1/2$ and $G_{\mathrm{poly}}(t) < 0$ for $|t| > 1/2$.
    Asymptotically $G_{\mathrm{poly}}(t) \sim -2/t^2$ as $t \to +\infty$.
  \item[\rm(ii)] From the Stirling expansion of $\psi^{(1)}$: $G_\Gamma(t) \sim +1/(2t^2)$ as
    $t \to +\infty$.
  \item[\rm(iii)] The following identity holds (from the reflection formula for $\psi^{(1)}$):
    \[
      \Re\!\left[\psi^{(1)}\!\left(\tfrac{1}{4}+\tfrac{it}{2}\right)\right]
      + \Re\!\left[\psi^{(1)}\!\left(\tfrac{3}{4}+\tfrac{it}{2}\right)\right]
      = 2\pi^2\operatorname{sech}^2(\pi t).
    \]
  \item[\rm(iv)] $G_{\mathrm{ana}}(0) = 8 - \tfrac{1}{4}\psi^{(1)}(1/4) \approx 8 - 4.30 = 3.70 > 0$.
  \item[\rm(v)] $G_{\mathrm{ana}}(t) \sim -3/(2t^2) < 0$ for $t \to +\infty$.
\end{enumerate}
In particular, $G_{\mathrm{ana}}$ changes sign: there exists $t_* > 0$ such that
$G_{\mathrm{ana}}(t_*) = 0$, $G_{\mathrm{ana}}(t) > 0$ for $t < t_*$, and $G_{\mathrm{ana}}(t) < 0$ for $t > t_*$.
\end{proposition}

\begin{proof}
(i): Direct from~\eqref{eq:gpoly}.
(ii): The Stirling expansion $\psi^{(1)}(s) \approx s^{-1} + (2s^2)^{-1} + \cdots$ for $|s| \to \infty$
applied to $s = 1/4+it/2 \approx it/2$ gives
$\psi^{(1)}(1/4+it/2) \approx -2i/t - 2/t^2 + O(t^{-3})$,
so $\Re\psi^{(1)} \approx -2/t^2$ and $G_\Gamma \approx 1/(2t^2) > 0$.
(iii): Reflection formula $\psi^{(1)}(z)+\psi^{(1)}(1-z) = \pi^2/\sin^2(\pi z)$ applied to
$z = 1/4+it/2$ and the duplication formula (Abramowitz--Stegun \cite{AS1964} \S6.4).
(iv): Numerical evaluation using $\psi^{(1)}(1/4) = \pi^2 + 8G \approx 17.20$ where $G$ is
Catalan's constant.
(v): Combining (i) and (ii): $G_{\mathrm{ana}} \approx -2/t^2 + 1/(2t^2) = -3/(2t^2)$.
\end{proof}

\subsection{The arithmetic obstruction}

\begin{proposition}[The wall]\label{prop:wall}
The inequality $G_\zeta(t) > -G_{\mathrm{ana}}(t)$ for all $t > t_*$ is equivalent to RH.
More precisely:
\begin{enumerate}
  \item[\rm(i)] If $\rho_0 = 1/2 + b_0 + i\gamma_0$ with $b_0 \neq 0$ is a zero of $\zeta$, then
    $G(t) < 0$ for $t$ near $\gamma_0$, violating the curvature criterion.
  \item[\rm(ii)] Controlling $G_\zeta(t)$ in the critical line $\sigma = 1/2$ requires control
    over the distribution of zeros of $\zeta$---precisely the content of RH.
\end{enumerate}
\end{proposition}

\begin{proof}
For (i): if $\rho_0 = 1/2+b_0+i\gamma_0$ is off-line, the quartet
$\{\rho_0, 1-\rho_0, \bar\rho_0, \overline{1-\rho_0}\}$ contributes to $G(t)$ near $t = \gamma_0$
terms of the form $-b_0^2/((t-\gamma_0)^4)$ with negative leading coefficient, so
$G(t) \to -\infty$ as $t \to \gamma_0$ from outside the zero.  Hence $G < 0$ near $\gamma_0$.

For (ii): the Euler product $\zeta(s) = \prod_p(1-p^{-s})^{-1}$ converges absolutely only
for $\Re(s) > 1$.  On the critical line $\sigma = 1/2$, the series
$-(\zeta'/\zeta)(1/2+it) = \sum_n \Lambda(n)n^{-1/2-it}$ (von Mangoldt) diverges absolutely.
Conditional convergence depends on cancellation between prime powers, which is controlled
by the zero distribution of $\zeta$---i.e., by RH.
\end{proof}

%------------------------------------------------------------------
\section{The wall: precise statement}
\label{sec:wall}
%------------------------------------------------------------------

Combining the results above, the following is the minimal precise statement of the
remaining obstruction:

\begin{theorem}[The wall]\label{thm:wall}
The following are equivalent:
\begin{enumerate}
  \item[\rm(1)] RH.
  \item[\rm(2)] $G(t) = G_{\mathrm{ana}}(t) + G_\zeta(t) > 0$ for all $t \notin \{\gamma_\rho\}$.
  \item[\rm(3)] $G_\zeta(t) > -G_{\mathrm{ana}}(t) \sim 3/(2t^2)$ for all $t > t_*$.
  \item[\rm(4)] $f(t) = \xi(1/2+it) \in \LP$.
\end{enumerate}
The analytic sector $G_{\mathrm{ana}}(t)$ is computable from the functional equation alone.
Controlling $G_\zeta(t)$ unconditionally on $\sigma = 1/2$ is equivalent to~\rm{(1)}.
A proof of the Riemann Hypothesis requires either a direct proof of~\rm{(3)} or a proof
of~\rm{(4)} by methods that bypass the Euler product on the critical line.
\end{theorem}

%------------------------------------------------------------------
\section{Conclusions}
\label{sec:conclusions}
%------------------------------------------------------------------

We have derived the curvature formula~\eqref{eq:curvature-cartesian} as an explicit
expression for $\partial^2_{\sigma\sigma}\log|\xi|$ and shown it to be positive everywhere
outside zeros if and only if RH holds (Theorems~\ref{thm:rh-curvature}).  We have
further shown that the four frameworks (A)--(D) in the literature are all equivalent
to this single criterion (Theorem~\ref{thm:lp}).

The sector decomposition (Theorem~\ref{thm:sectors}) shows that the analytic sector
$G_{\mathrm{ana}} = G_{\mathrm{poly}} + G_\Gamma$ is explicitly computable and changes sign
at a threshold $t_*$; for $t > t_*$ the arithmetic sector $G_\zeta$ must compensate.
This compensation is provably equivalent to RH (Proposition~\ref{prop:wall}), and
constitutes the irreducible obstruction.

\medskip

\noindent\textbf{What is new.}
\begin{itemize}
  \item The explicit Hadamard formula for $\mathcal{C}(\gamma)$ in Cartesian coordinates~\eqref{eq:curvature-cartesian}.
  \item The resolution of the concavity direction: $\sigma$-convexity (not concavity) holds under RH.
  \item The exact sector decomposition with polygamma formulas (Theorem~\ref{thm:sectors},
    Proposition~\ref{prop:gana}).
  \item The unification Theorem~\ref{thm:lp} synthesizing frameworks (A)--(D) into one object.
\end{itemize}

\noindent\textbf{What is not new.}
The equivalences $\mathrm{RH}\iff\LP\iff\Lambda=0$ are due to de~Bruijn~\cite{deBruijn1950}
and Newman~\cite{Newman1976}.  The inequality $\Lambda\ge0$ is Rodgers--Tao~\cite{RodgersTao2020}.
The APS connection uses standard spectral geometry.  The $K$-theoretic formulation
follows Connes~\cite{Connes1999}.

%------------------------------------------------------------------
\begin{thebibliography}{99}

\bibitem{APS1975}
M.~F. Atiyah, V.~K. Patodi, and I.~M. Singer,
\textit{Spectral asymmetry and Riemannian geometry. I},
Math.\ Proc.\ Cambridge Philos.\ Soc.\ \textbf{77} (1975), 43--69.

\bibitem{AS1964}
M.~Abramowitz and I.~A. Stegun (eds.),
\textit{Handbook of Mathematical Functions},
National Bureau of Standards, Washington, D.C., 1964.

\bibitem{CNV1988}
G.~Csordas, R.~S. Norfolk, and R.~S. Varga,
\textit{The Riemann hypothesis and the Turán inequalities},
Trans.\ Amer.\ Math.\ Soc.\ \textbf{296} (1986), 521--541.

\bibitem{Connes1999}
A.~Connes,
\textit{Trace formula in noncommutative geometry and the zeros of the Riemann
  zeta function},
Selecta Math.\ (N.S.) \textbf{5} (1999), 29--106.

\bibitem{deBruijn1950}
N.~G. de~Bruijn,
\textit{The roots of trigonometric integrals},
Duke Math.\ J.\ \textbf{17} (1950), 197--226.

\bibitem{Newman1976}
C.~M. Newman,
\textit{Fourier transforms with only real zeros},
Proc.\ Amer.\ Math.\ Soc.\ \textbf{61} (1976), 245--251.

\bibitem{RodgersTao2020}
B.~Rodgers and T.~Tao,
\textit{The de Bruijn--Newman constant is non-negative},
Forum Math.\ Pi \textbf{8} (2020), e6.

\end{thebibliography}

\end{document}
